\section{Instrument Validation}

\label{sec:instrumentos}

Before evaluating the system against real papers (section~\ref{sec:experimentos}), this section validates each instrument separately against known ground truth: D3 against calibration pairs with human judgment, D2 against the evaluation set, and D1 C$_F$ against the coverage of the same set.

\subsection{D3: Calibration Ladder}
\label{ssec:d3cal}

The Jaccard distance was validated on five theorem pairs for which a human judge (the first author) determined the expected relationship between their proofs and signed each label before running the measurement. The pairs cover the spectrum: same proof (\textit{self} control), genuinely different proofs, apparently different proofs that turned out identical after formalization, and unrelated theorems (cross control).

\begin{table}[htbp]
\centering
\small
\setlength{\tabcolsep}{5pt}
\begin{tabular}{llrrrr}
\toprule
Pair & Judge label & Jaccard & Distance & $|\cap|$ & $|\cup|$ \\
\midrule
T08a vs.\ T08a & control: same proof      & 1{,}0000 & 0{,}0000 & 9 & 9 \\
T07a vs.\ T07b & \textit{same disguised}    & 0{,}5000 & 0{,}5000 & 1 & 2 \\
T08a vs.\ T08b & \textit{genuinely different} & 0{,}2778 & 0{,}7222 & 5 & 18 \\
T09a vs.\ T09b & \textit{genuinely different} & 0{,}0000 & 1{,}0000 & 0 & 6 \\
T07 vs.\ T08   & control: unrelated   & 0{,}0000 & 1{,}0000 & 0 & 10 \\
\bottomrule
\end{tabular}
\caption{D3 calibration over five pairs with prior human labeling. Premise sets are post-filter and deduplicated.}
\label{tab:d3cal}
\end{table}

\textbf{T08 (distance 0{,}7222).} Two genuinely different proofs of the irrationality of $\sqrt{2}$ yield high distance. The five shared premises are foundational arithmetic lemmas; the remaining ones, exclusive to each side, capture the divergent strategies (divisibility by primes versus 2-adic valuation). The threshold $\theta = 0{,}5$ correctly classifies the pair as distant proofs.

\textbf{T09 (distance 1{,}0).} The two proofs of the Gauss sum, induction with \texttt{sum\_range\_succ} versus closed formula with \texttt{Finset.sum\_range\_id}, share no premises after filtering. Each side contributes six premises exclusive to its strategy, and the maximum distance correctly reflects the difference labeled by the judge.

\textbf{T07 (distance 0{,}50): degenerate point.} The two proofs of the infinitude of primes (Euclid and Euler) fall exactly on the threshold. After formalization, both collapsed to the same invocation of \texttt{Nat.exists\_infinite\_primes}. With only two total premises after filtering and one shared, the Jaccard is fragile: minor changes in filter behavior would invert the verdict. The judge had labeled them as \textit{same disguised}: distinct statements in the LaTeX source, identical in the elaborated proof term. This pair is not informative for calibration.

\textbf{Negative control.} The distance between a number theory theorem and an analysis theorem is 1{,}0 with empty intersection, confirming that D3 does not produce spurious similarity across distinct domains.

\textbf{What this ladder calibrates.} The five pairs are consistent with human judgment under $\theta = 0{,}5$. However, three of them (T09 and the two controls) fall at the extremes of the range and do not constrain the choice of threshold: any value in $(0,1)$ classifies them correctly. The only effective bounds come from T08, which requires $\theta < 0{,}7222$, and from T07, which requires $\theta \geq 0{,}50$. Since T07 is degenerate (union of two premises, intersection of one, such that a minor change in filter behavior would invert the verdict), the lower bound is not reliable. The admissible interval supported by evidence thus reduces to $\theta < 0{,}7222$, and the value $\theta = 0{,}5$ is a design choice within that range, not a calibrated optimum. Narrowing that interval requires pairs whose distance falls in the interior of the range and whose premise sets are large enough for the Jaccard ratio to be stable; constructing them is future work (section~\ref{sec:conclusion}).

\textbf{D3 measures formalizations, not ideas.} The T07 case exposes a limit of the metric that goes beyond the choice of threshold. Mathlib typically contains a single canonical proof of each classical theorem, so that two strategies a mathematician would consider distinct can converge to the same lemma upon being formalized. The consequence is that D3 measures distance between \textit{formalizations}, and the relationship between that distance and the distance between the underlying informal ideas remains unestablished: a distance of 1{,}0 may be due to the formalizer choosing different lemmas, and a low distance to convergence to the same Mathlib lemma. We report D3 as relative to the formalization over Mathlib v4.29.0. Quantifying how often that convergence occurs requires a systematic experiment across multiple theorems and multiple formalizers, which remains as future work.

\subsection{D2: Triviality Filter}

D2 was evaluated on 24 of the 26 theorems in the evaluation set. Two cases (T20 and T21) were not formalized in this run for reasons unrelated to the instrument. The 24 evaluated cover trivial-by-design cases, classics present in Mathlib, pairs with different proofs, near statements, LLM-generated cases, and deliberate failure cases.

\begin{table}[htbp]
\centering
\small
\setlength{\tabcolsep}{5pt}
\begin{tabular}{llll}
\toprule
Case & Statement & Tactic & Expectation \\
\midrule
T14 & sum of four evens is even        & \texttt{aesop}    & trivial \\
T15 & $2 + 2 = 4$                        & \texttt{decide}   & trivial \\
T16 & $n + 0 = n$                        & \texttt{norm\_num} & trivial \\
T17 & $n \leq n + 1$                     & \texttt{norm\_num} & trivial \\
T19 & LLM-generated about evens       & \texttt{aesop}    & ambiguous \\
T22 & $n$ even $\Rightarrow n+0$ even      & \texttt{norm\_num} & designed for D1 \\
\bottomrule
\end{tabular}
\caption{The six statements that D2 closed with a tactic from \texttt{T\_AUTO}. The first four had a defined triviality expectation; T19 and T22 are discussed in the text.}
\label{tab:d2trivial}
\end{table}

T19 was generated by an LLM instructed to state and prove an original theorem about even numbers; D2 closed it with \texttt{aesop}, which is consistent with the hypothesis that models tend to produce trivial statements, but its expectation in the evaluation set was not binary. T22 is logically equivalent to ``if $n$ is even then $n$ is even''; D2 closed it correctly, but the case was constructed to test syntactic equivalence in D1, not triviality. Neither admits a correct/incorrect classification for D2.

\begin{table}[htbp]
\centering
\small
\setlength{\tabcolsep}{5pt}
\begin{tabular}{p{4.6cm}p{6.4cm}l}
\toprule
Cases & Category & D2 Result \\
\midrule
T01--T06                & classics present in Mathlib            & non-trivial \\
T07a, T08a, T09a        & pairs with different proofs          & non-trivial \\
T10--T13                & near statements (odd primes, AM--GM) & non-trivial \\
T18                     & sum of first $n$ odds $= n^2$ & non-trivial \\
T23                     & connected acyclic graph is a tree         & non-trivial \\
T24                     & coherent sheaves on Noetherian scheme & non-trivial \\
T25                     & $n$ even $\iff 2 \mid n$                  & non-trivial \\
T26                     & sum of $n$ evens is even                 & non-trivial \\
\bottomrule
\end{tabular}
\caption{The 18 statements that D2 correctly classified as non-trivial.}
\label{tab:d2notrivial}
\end{table}

\textbf{Notable cases.} T18 is a control trap: it requires induction and D2 correctly did not close it. T23 was also not closed in this run, unlike previous runs where \texttt{tauto} closed it on a conjunctive definition; the Lean statement effectively used here differs from the one that produced that false positive. T14, in contrast, was closed by \texttt{aesop} in this run, whereas in previous runs with a more complex statement the same tactic exceeded the budget by more than an order of magnitude. Both cases illustrate D2's sensitivity to the concrete statement produced by the formalizer, discussed in section~\ref{sec:limitaciones}.

\subsubsection*{Aggregate}

Over the 24 evaluated theorems, D2 succeeds on 22 (91{,}7\,\%). T19 and T22 remain in the denominator but not in the numerator, lacking a binary expectation against which to measure success; over the 22 cases with defined expectation no false positives or false negatives are recorded.

\subsection{D1 C$_F$: Formal Corpus Coverage}

The C$_F$ branch found a match for the 18 non-trivial theorems that reached that stage. The remaining six do not register a match because D1 was never executed on them: they are exactly the six that D2 detected as trivial, on which the tree applied short-circuit. The absence of match in those cases reflects non-execution, not a search failure.

\textbf{This result should not be read as measured coverage.} The C$_F$ match criterion accepts the first result returned by Leandex without applying a similarity threshold, because the current API version does not provide scores (section~\ref{sec:pipeline}). Such a criterion can only return absence of match when the API responds empty, so that 18 matches out of 18 queries does not distinguish real coverage from absence of filtering. The 18 cases correspond to classical theorems that are indeed in Mathlib, and the composition of the evaluation set (dominated by canonical results) makes that outcome expected under any criterion. Establishing the precision of C$_F$ requires a similarity threshold or a manual verification of each match, and neither is available in this version.

The C$_I$ branch was not activated for any theorem. The cause is not the embedding threshold but the entry condition: C$_I$ is only consulted when D2 is negative and C$_F$ found no match, and that combination never occurred. All non-trivial cases had a C$_F$ match, and all those that did not were stopped at D2. The evaluation set, by composition, does not exercise that branch.

\subsection{Summary}

\begin{table}[htbp]
\centering
\small
\setlength{\tabcolsep}{4pt}
\renewcommand{\arraystretch}{1.15}
\begin{tabular}{p{1.9cm}p{2.6cm}p{6.4cm}p{3.4cm}}
\toprule
Instrument & Ground truth & Result & Status \\
\midrule
D3 &
5 pairs with prior human label &
5/5 consistent with the judge under $\theta = 0{,}5$. Only T08 and T07 constrain the threshold, and T07 is degenerate ($|\cup| = 2$): the evidence supports $\theta < 0{,}7222$, not a point value. &
Functional; $\theta$ not calibrated \\
\addlinespace
D2 &
24 theorems from the evaluation set &
22/24 successes (91{,}7\,\%); T19 and T22 in the denominator without binary expectation. No false positives or negatives over the remaining 22. &
Functional \\
\addlinespace
D1 C$_F$ &
18 non-trivial theorems &
18 matches over 18 queries, but without a similarity threshold the result does not distinguish coverage from absence of filtering. &
Inconclusive \\
\addlinespace
D1 C$_I$ &
--- &
Branch not reached: the entry condition was never met with this set. &
Not evaluated \\
\bottomrule
\end{tabular}
\caption{Validation status of each instrument. Only D2 admits an interpretable success measure with the current evaluation set.}
\label{tab:instrumentos}
\end{table}